% Einheit 2 von 3 – Binomische Formeln (bank/binomische-formeln/e2.jsonl, zusammenbau v0.1)
%% TODO Zweigzeile Teil 1: was hier gelernt wird (Satz in Schülersprache aus dem Einheitstitel) – Einheit 2
\einheitenkopf[e2]{Einheit 2 von 3 · Binomische Formeln}
\zweigzeile{ab Kl. 8 · P10}
\setcounter{aufgabe}{12}

%% TODO Titel: Ich-kann-Satz fehlt in der Bank (Platzhalter: Kettenname) – Nr. 13
%% TODO Anweisung: ein Satz über der ersten Teilaufgabe fehlt in der Bank – Nr. 13
\begin{aufgabe}{Gleiche Klammer zweimal?}
\begin{teile}
\teil $(x + 9)^2$ – Quadrat einer Klammer? Wenn ja: schreib es als Klammer mal Klammer. \janein
\end{teile}
\end{aufgabe}

%% TODO Titel: Ich-kann-Satz fehlt in der Bank (Platzhalter: Kettenname) – Nr. 14
%% TODO Anweisung: ein Satz über der ersten Teilaufgabe fehlt in der Bank – Nr. 14
\begin{aufgabe}{Binomische Formeln}
%% TODO \rechenplatz{n}: Schrittzahl des Grundfalls fehlt in der Bank (nur bei fester Schreibform, 2.3 b)
\begin{teile}
\teil Welche binomische Formel passt, und was ist $a$, was ist $b$ in der Formel? $(x + 14)^2$ \leerfeld Formel; a = \leerfeld , b = \leerfeld
\end{teile}
\begin{gleichungsraster}[2]
\gl{\text{Multipliziere aus: $(x + 4)^2$}} &
\gl{\text{Multipliziere aus: $(x - 9)^2$}} \\
\gl{\text{Multipliziere aus: $(x + 8) \cdot (x - 8)$}} &
\gl{\text{Welche Formel passt? Multipliziere aus: $(x + 9) \cdot (x - 9)$}} \\
\gl{\text{Multipliziere aus: $(3x + 1)^2$}} &
\gl{\text{Multipliziere aus: $(x + 3y)^2$}} \\
\gl{\text{Multipliziere aus: $2 \cdot (x + 4)^2$}} &
\gl{\text{Multipliziere aus: $-(x + 9)^2$}} \\
\gl{\text{Multipliziere aus und fasse zusammen: $(x + 1)^2 - 6$}} &
\end{gleichungsraster}
\begin{teile}
\teil Die Parabel $p$ hat die Gleichung $y = (x + 4)^2 - 7$. Weise nach, dass $y = x^2 + 8x + 9$ dieselbe Parabel beschreibt. \hfill (P10 2017 OS)
\end{teile}
\end{aufgabe}

%% TODO Titel: Ich-kann-Satz fehlt in der Bank (Platzhalter: Kettenname) – Nr. 15
%% TODO Anweisung: ein Satz über der ersten Teilaufgabe fehlt in der Bank – Nr. 15
\begin{aufgabe}{Kopfrechnen mit der Formel (Vorrat)}
\begin{teile}
\teil $41^2$ – mit einer binomischen Formel im Kopf? \leerfeld
\end{teile}
\end{aufgabe}

%% TODO Titel: Ich-kann-Satz fehlt in der Bank (Platzhalter: Kettenname) – Nr. 16
%% TODO Anweisung: ein Satz über der ersten Teilaufgabe fehlt in der Bank – Nr. 16
\begin{aufgabe}{Binomische Formeln – Fehler finden, Begründen}
\begin{teile}
\teil Finde den Fehler und rechne richtig: \rechnung{(y + 9)^2 &= y^2 + 81}
\teil Tim meint: $(x + 4)^2 = x^2 + 16$. Prüfe mit $x = 2$ und begründe, warum das nicht stimmt.
\end{teile}
\end{aufgabe}
